% ---------------------------------------------------------------------------
% Theorem environments
% ---------------------------------------------------------------------------
\theoremstyle{plain}
\newtheorem{theorem}{Theorem}[section]
\newtheorem{lemma}[theorem]{Lemma}
\newtheorem{corollary}[theorem]{Corollary}
\newtheorem{proposition}[theorem]{Proposition}
\newtheorem{conjecture}[theorem]{Conjecture}
\newtheorem{question}[theorem]{Question}
\newtheorem{problem}[theorem]{Problem}
\theoremstyle{definition}
\newtheorem{definition}[theorem]{Definition}
\newtheorem{remark}[theorem]{Remark}
\newtheorem{example}[theorem]{Example}

% Keep a theorem heading, formal-status line and opening statement together.
\AtBeginEnvironment{theorem}{\Needspace{8\baselineskip}}
\AtBeginEnvironment{lemma}{\Needspace{8\baselineskip}}
\AtBeginEnvironment{corollary}{\Needspace{8\baselineskip}}
\AtBeginEnvironment{proposition}{\Needspace{8\baselineskip}}
\AtBeginEnvironment{example}{\Needspace{8\baselineskip}}

% The Lean proof is the proof; the text gives a faithful outline of it.
\newenvironment{outline}[1][Outline of the formal proof]{\begin{proof}[#1]}{\end{proof}}
% Results that are not formalized get complete proofs, labelled as such.
\newenvironment{informalproof}{\begin{proof}[Proof (not formalized)]}{\end{proof}}

% ---------------------------------------------------------------------------
% Lean references
%   \lean{name}            a declaration, relative to the namespace BoundedArityDecoders
%   \leanname{name}        typeset an upstream or qualified Lean identifier without indexing it
%   \formal{file}{names}   the formal counterpart of a numbered result: a file relative to
%                          lean/BoundedArityDecoders/ and a comma-separated list of declarations
%   \partformal{file}{names}{gap}
%                          a numbered result whose formal counterpart covers all but the
%                          part described in {gap}; the gap is proved or cited in the text
%   \notformalized         a numbered result that has no formal counterpart
%   \cited{where}          a result of an earlier paper used without proof: where it is proved,
%                          and whether it is Lean-verified there
% ---------------------------------------------------------------------------
% Lean names are typeset like paths, with line breaks allowed after "_", ".", "/", "-" and camelCase transitions.
\emergencystretch=1.5em
\DeclareUrlCommand\leannameraw{%
  \def\UrlFont{\ttfamily\upshape}%
  \Urlmuskip=0mu\relax
  \def\UrlBreakPenalty{500}%
  \def\UrlBreaks{\do\_\do\.\do\/\do\-\do\a\do\b\do\c\do\d\do\e\do\f\do\g\do\h\do\i\do\j\do\k\do\l\do\m\do\n\do\o\do\p\do\q\do\r\do\s\do\t\do\u\do\v\do\w\do\x\do\y\do\z}%
}
\renewcommand{\UrlBreaks}{\do\_\do\.\do\/\do\-}
\newcommand{\leanname}[1]{\ifmmode\text{\leannameraw{#1}}\else\leannameraw{#1}\fi}
\newcommand{\lean}[1]{\leanname{#1}}
\newcommand{\leanfile}[1]{\leanname{#1}}
\newcommand{\leanlist}[1]{\begingroup\def\leansep{}%
  \renewcommand*{\do}[1]{\leansep\leannameraw{##1}\def\leansep{, }}\docsvlist{#1}\endgroup}
\newcommand{\formal}[2]{\par\noindent{\small\upshape\raggedright\textsf{Lean:} \leanlist{#2}\unskip
  \nobreak\hfil\penalty50\hskip1em\hbox{}\nobreak\hfill[\leanfile{#1}]\parfillskip=0pt\par}%
  \smallskip\noindent\ignorespaces}
\newcommand{\partformal}[3]{\par\noindent{\small\upshape\raggedright\textsf{Lean (partial):} \leanlist{#2}\unskip
  \nobreak\hfil\penalty50\hskip1em\hbox{}\nobreak\hfill[\leanfile{#1}]\parfillskip=0pt\par}%
  \par\noindent{\small\upshape\textsf{Not formalized:} #3}\par\smallskip\noindent\ignorespaces}
\newcommand{\notformalized}{\par\noindent{\small\upshape\textsf{Not formalized;} a complete proof is given
  below.}\par\smallskip\noindent\ignorespaces}
\newcommand{\cited}[1]{\par\noindent{\small\upshape\textsf{Cited:} #1}\par\smallskip\noindent\ignorespaces}

% ---------------------------------------------------------------------------
% Notation
% ---------------------------------------------------------------------------
\newcommand{\N}{\mathbb{N}}
\newcommand{\Z}{\mathbb{Z}}
\newcommand{\Q}{\mathbb{Q}}
\newcommand{\R}{\mathbb{R}}
\newcommand{\C}{\mathbb{C}}
\newcommand{\F}{\mathbb{F}}
\newcommand{\T}{\mathsf{T}}
\newcommand{\clog}[1]{\lceil \log_2 #1 \rceil}
\newcommand{\abs}[1]{\lvert #1 \rvert}
\newcommand{\card}[1]{\lvert #1 \rvert}
\newcommand{\eps}{\varepsilon}
\newcommand{\TCP}{\mathrm{TCP}}
\newcommand{\TM}{\mathrm{TM}}
\newcommand{\rot}{\operatorname{rot}}
\newcommand{\Mat}{\operatorname{Mat}}
\newcommand{\End}{\operatorname{End}}
\newcommand{\Hom}{\operatorname{Hom}}
\newcommand{\Rec}{\operatorname{Rec}}
\newcommand{\sw}{\operatorname{sw}}
\newcommand{\diam}{\operatorname{diam}}
\newcommand{\OSR}{\operatorname{rk}_{\otimes}}
\DeclareMathOperator{\rank}{rank}
\DeclareMathOperator{\spn}{span}
\DeclareMathOperator{\poly}{poly}
